\begin{abstract}
Evolution strategies (ES) can distribute model evaluations by exchanging only scalar
fitness scores, but each device then usually reconstructs the same parameter update. We
compare this replication with splitting the reconstruction across devices and combining
partial updates with an all-reduce. The faster choice depends on whether the computation
saved outweighs the communication added. Using one JAX implementation, we measure this
tradeoff for full-rank and low-rank perturbations on A100 GPUs and TPU v5e chips. On a
transformer-block benchmark with eight devices in one host, splitting gives
1.5--5.6$\times$ speedups for the unpaired full-rank variants; for low-rank perturbations at the
larger matrix width, it takes 4--46\% longer than replication. The placement preference
carries over to Qwen2.5-0.5B on the TPU, while a slow connection between two A100 hosts
reverses the full-rank preference for large updates. Supporting learning experiments find
similar directional alignment and observed final Countdown accuracy across the tested
perturbations, with full rank ahead at the first training checkpoint. These checks do not
establish equivalent learning efficiency. Code, configurations, and results are released.
\end{abstract}
